% from main.tex
%%%%%%%%%%  scan idx 687  |  folio 674  %%%%%%%%%%
régulier ou non). Il a par contre le désavantage (4.6) sur $A(X)$, $\Gtop^\bullet(X)$ de ne pas avoir tel quel un caractère covariant par rapport aux morphismes propres (même pour des schémas lisses quasi-projectifs sur un corps)~; il n'acquiert ce caractère qu'une fois tensorisé par $\Q$.

Autre avantage de $A(X)$ sur ses concurrents $\Gtop^\bullet(X)$ et $\operatorname{Gr}^\bullet(X)$, c'est qu'on dispose d'un homomorphisme de $A(X)$ dans $H^{2\bullet}(X,\Zell(\bullet))$, permettant d'interpréter les classes de Chern cohomologiques comme images des classes de Chern à valeurs dans $A(X)$. Cet homomorphisme en général ne se factorise pas par $\Gtop^\bullet(X)$ (4.7), et il est pour le moins douteux qu'il existe en général un homomorphisme de $\operatorname{Gr}^\bullet(X)$ dans $H^{2\bullet}(X,\Zell(\bullet))$ qui transforme classes de Chern en classes de Chern (Cf. no 5). Bien entendu, ces difficultés disparaissent lorsqu'on tensorise par $\Q$, et on trouve un homomorphisme
\begin{equation}\tag{4.2}
 \Gtop^\bullet(X)\longrightarrow H^{2\bullet}(X,\Qell(\bullet))=
 H^{2\bullet}(X,\Zell(\bullet))\otimes_{\Zell}\Qell
\end{equation}
qui transforme classes de Chern en classes de Chern.

\textbf{4.2.} Pour établir que l'homomorphisme surjectif $\varphi$ de (4.1) est un isomorphisme modulo torsion, on utilise l'homomorphisme de Chern (déduit de la théorie des classes de Chern [9])
\[
 K^\bullet(X)\longrightarrow \Ch(A(X)),
\]
(où le deuxième membre est l'anneau de Chern de $A(X)$ défini dans Exp.~V), qui est compatible comme on voit aisément avec la filtration topologique de $K^\bullet(X)$ et la filtration naturelle de $\Ch(A(X))$, et donne par passage aux gradués associés un homomorphisme de groupes gradués
\begin{equation}\tag{4.3}
 \psi:\Gtop^\bullet(X)\longrightarrow A(X).
\end{equation}
Ceci posé, on va prouver les formules
\begin{equation}\tag{4.4}
\left\{
\begin{aligned}
 \psi\varphi(x)&=(-1)^{i-1}(i-1)!\,x\quad \bmod \text{ torsion}\quad \text{pour } x\in A^i(X),\\
 \varphi\psi(x)&=(-1)^{i-1}(i-1)!\,x\quad \bmod \text{ torsion}\quad \text{pour } x\in \Gtop^i(X);
\end{aligned}
\right.
\end{equation}
dont la première implique d'ailleurs la seconde, grâce au fait que $\varphi$ est surjectif. Ces formules impliquent évidemment que $\varphi$ et $\psi$ sont des isomorphismes modulo torsion, et qu'à un facteur $(-1)^{i-1}(i-1)!$ près en degré $i$, ils sont inverses l'un de l'autre.

Pour prouver la première formule (4.4), on peut évidemment supposer que $x$ est défini par un sous-schéma fermé irréductible $Y$ de $X$ de codimension $i$, auquel cas la formule équivaut à
\begin{equation}\tag{4.5}
 c_i(\mathcal O_Y)=(-1)^{i-1}(i-1)!\,\cl(Y)\quad \bmod
